\documentclass[journal]{IEEEtran}
\usepackage[T1]{fontenc}
\usepackage{amsmath}
\title{A Simple Sensor Averaging Method}
\author{Ada Author}
\begin{document}
\maketitle
\begin{abstract}
This fictional example illustrates an arithmetic mean of sensor readings.
It is demonstration content, not a scientific finding.
\end{abstract}
\section{Introduction}
We summarize readings using a simple mean, following fictional guides
by Alpha~\cite{guide} and Beta~\cite{method}.
\section{Method}\label{sec:method}
We compute $\bar{x}=N^{-1}\sum_{i=1}^{N}x_i$.
Here N is the positive number of readings included in one batch.
N may differ between batches; we compute a separate mean for each batch.
All readings in a batch measure the same quantity in the same units.
\section{Results}\label{sec:results}
The example temperature readings are 10, 12, and 14 degrees Celsius.
Their mean is $(10+12+14)/3=12$ degrees Celsius.
\section{Limitations}\label{sec:limitations}
An outlier can shift the arithmetic mean. This example does not claim
robustness to outliers, and no robust estimator is evaluated.
We follow Delta's fictional limitations checklist~\cite{limits}.
\section{Conclusion}
The mean provides a compact summary of these example readings.
\bibliographystyle{IEEEtran}
\bibliography{paper}
\end{document}
